\documentclass[12pt,a4paper]{exam}

% Paquetes necesarios
\usepackage[utf8]{inputenc}
\usepackage[T1]{fontenc}
\usepackage[spanish]{babel}
\usepackage{siunitx} % Para unidades
\usepackage{multicol}
\usepackage{geometry}

% Márgenes muy pequeños para 4 exámenes en 1 hoja A4
% top, bottom, left, right
\geometry{top=1cm, bottom=1cm, left=1cm, right=1cm, columnsep=1cm}

\pagestyle{empty} % sin encabezados ni pies

% Definir una regla para el nombre (con menos espacio)
\newcommand{\nombre}{\noindent\makebox[\linewidth]{\textbf{Nombre y Apellido:}\enspace\hrulefill}\vspace{0.25cm}}

% Nota sobre Kelvin (más corta)
\newcommand{\notakelvin}{\small \textit{(Recuerde usar Kelvin para cálculos de $\eta$ ideal.)}}

\begin{document}

% Usamos un solo entorno multicols para toda la página
% Col 1: Tema A (arriba), Tema C (abajo)
% Col 2: Tema B (arriba), Tema D (abajo)
\begin{multicols}{2}

% ================= TEMA A ==================
\nombre
\begin{center}
    \textbf{RECUPERATORIO - TEMA A} \\
    \notakelvin
\end{center}
\vspace{0.1cm}
\begin{questions}
 \question (Carnot) Una máquina ideal opera entre $T_1 = \SI{600}{\celsius}$ y $T_2 = \SI{30}{\celsius}$. Absorbe $\SI{2500}{J}$ ($Q_1$) de la fuente caliente.
    \begin{parts}
        \part[2] Calcule la eficiencia térmica ideal ($\eta_{ideal}$).
        \part[2] Determine el trabajo neto ($W$) realizado.
        \part[1] Calcule el calor desechado ($Q_2$).
    \end{parts}    
 \question (Máquina real) Una máquina térmica realiza un trabajo neto $W = \SI{500}{J}$ y desecha $Q_2 = \SI{800}{J}$ al sumidero frío en cada ciclo.
      \begin{parts}
        \part[2] Calcule el calor absorbido ($Q_1$).
        \part[3] Calcule la eficiencia térmica real ($\eta_{real}$) de la máquina.
    \end{parts}
\end{questions}
\vfill % Empuja el siguiente examen hacia abajo

% ================= TEMA C ==================
\nombre
\begin{center}
    \textbf{RECUPERATORIO - TEMA C} \\
    \notakelvin
\end{center}
\vspace{0.1cm}
\begin{questions}
 \question (Carnot) Una máquina ideal opera entre $T_1 = \SI{650}{\celsius}$ y $T_2 = \SI{20}{\celsius}$. Absorbe $\SI{2000}{J}$ ($Q_1$) de la fuente caliente.
    \begin{parts}
        \part[2] Calcule la eficiencia térmica ideal ($\eta_{ideal}$).
        \part[2] Determine el trabajo neto ($W$) realizado.
        \part[1] Calcule el calor desechado ($Q_2$).
    \end{parts}    
 \question (Máquina real) Una máquina térmica realiza un trabajo neto $W = \SI{400}{J}$ y desecha $Q_2 = \SI{700}{J}$ al sumidero frío en cada ciclo.
      \begin{parts}
        \part[2] Calcule el calor absorbido ($Q_1$).
        \part[3] Calcule la eficiencia térmica real ($\eta_{real}$) de la máquina.
    \end{parts}
\end{questions}
\vfill % Termina la primera columna

\columnbreak % Pasa a la segunda columna

% ================= TEMA B ==================
\nombre
\begin{center}
    \textbf{RECUPERATORIO - TEMA B} \\
    \notakelvin
\end{center}
\vspace{0.1cm}
\begin{questions}
 \question (Carnot) Una máquina ideal opera entre $T_1 = \SI{550}{\celsius}$ y $T_2 = \SI{25}{\celsius}$. Absorbe $\SI{3000}{J}$ ($Q_1$) de la fuente caliente.
    \begin{parts}
        \part[2] Calcule la eficiencia térmica ideal ($\eta_{ideal}$).
        \part[2] Determine el trabajo neto ($W$) realizado.
        \part[1] Calcule el calor desechado ($Q_2$).
    \end{parts}    
 \question (Máquina real) Una máquina térmica realiza un trabajo neto $W = \SI{600}{J}$ y desecha $Q_2 = \SI{1000}{J}$ al sumidero frío en cada ciclo.
      \begin{parts}
        \part[2] Calcule el calor absorbido ($Q_1$).
        \part[3] Calcule la eficiencia térmica real ($\eta_{real}$) de la máquina.
    \end{parts}
\end{questions}
\vfill % Empuja el siguiente examen hacia abajo

% ================= TEMA D ==================
\nombre
\begin{center}
    \textbf{RECUPERATORIO - TEMA D} \\
    \notakelvin
\end{center}
\vspace{0.1cm}
\begin{questions}
 \question (Carnot) Una máquina ideal opera entre $T_1 = \SI{500}{\celsius}$ y $T_2 = \SI{15}{\celsius}$. Absorbe $\SI{3500}{J}$ ($Q_1$) de la fuente caliente.
    \begin{parts}
        \part[2] Calcule la eficiencia térmica ideal ($\eta_{ideal}$).
        \part[2] Determine el trabajo neto ($W$) realizado.
        \part[1] Calcule el calor desechado ($Q_2$).
    \end{parts}    
 \question (Máquina real) Una máquina térmica realiza un trabajo neto $W = \SI{750}{J}$ y desecha $Q_2 = \SI{1200}{J}$ al sumidero frío en cada ciclo.
      \begin{parts}
        \part[2] Calcule el calor absorbido ($Q_1$).
        \part[3] Calcule la eficiencia térmica real ($\eta_{real}$) de la máquina.
    \end{parts}
\end{questions}
\vfill % Termina la segunda columna

\end{multicols}

\end{document}